% Transcription — fonds Alexandre Grothendieck, shelfmark 81, batch 3 (pages 41-54).
% Established from the facsimile archives/batches/81/batch-03.pdf (a mirror of
% https://grothendieck.umontpellier.fr/81.pdf), per the skill
% transcribe-grothendieck. The batch file carries no cover sheet: PDF page N
% is archivists' page 40+N; the pencilled numbers bottom-left read « 41 »,
% « 42 », « 43 », « 44 », « 48 », « 49 », « 50 », « 51 », « 52 », « 53 »,
% « 54 » at the expected PDF pages (1-4, 8-14), and the others fall in
% sequence. He did not paginate these leaves.
%
% Pass: Opus 5.5 (claude-opus-5-5), 2026-09-23 — first pass, unchecked against the pages by a human.
%
% All fourteen pages carry a \page. Page 54 has no word on it (pencil
% sketches continuing page 53's gluing pictures) and is kept as a \page with
% a \note describing the drawings. Nothing is skipped as unrelated.
%
% What the batch is. Combinatorics of a finite « polygon »: a finite set S of
% cardinal N ≥ 3 with a « bi-circulation » {u, u^{-1}} (a cyclic order up to
% reversal), i.e. the vertices of a disc's boundary circle.
%   41-43  Segments of S (T = {s_0, u s_0, ..., u^{i-1} s_0}, 0 < i < N);
%          complements, intersections, unions; endpoints; the bijection
%          segments <-> S × [1, N-1] once an orientation is chosen; the
%          segments contained in a proper subset S' (its components, read
%          topologically on a polygon (X, S)); Proposition: for disjoint
%          A, B ≠ ∅, S the equivalent conditions a) disjoint enclosing
%          segments, b) no crossing pairs {a,a'}, {b,b'}, c)/c') each lies
%          in a component of the other's complement; « A, B non parallèles »;
%          margin d): convex hulls P_A ∩ P_B = ∅ for S on a convex curve.
%   44     « Structure à étudier » : partitions of S into mutually non
%          parallel parts (« partitions non parallèles »); intuition that
%          they reduce to degree-2 partitions (chord diagrams) on a smaller
%          bi-circulated set. A circle drawn below.
%   45-47  Families Σ of nested-or-disjoint segments (σ1, σ2), the « cores »
%          A* = A \ ∪_{B ⊊ A} B, their disjointness, the partition Σ̂* of S;
%          (σ3); page 47: whether Σ̂* determines Σ — no, a counterexample
%          with three chords a_1a_1', bb', a_2a_2' (drawn).
%   48-50  Brown ink. « Relation d'équivalence non parallèle » R on S;
%          R-saturated segments; Prop: minimal saturated segments are the
%          classes that are segments (« anglets »); Cor: ≥ 2 anglets unless
%          R is the coarse relation; the induction removing anglets
%          (R_1 = R|S_1 ...), the tower Σ_1, Σ_2, ... / T_1, T_2, ...
%   51-52  Black ink, a restart for a « quasi-segment » I: non parallel pairs
%          ε, ε' ∈ 𝔓_2(I) and parts; « Partitions non parallèles de I »
%          from a family Σ of segments with a), b), c); A*, Σ̂, Σ̂*; NB every
%          non parallel partition arises so, in several ways; Σ_max. Page 52
%          stops after four lines.
%   53-54  A different subject on a square sheet: gluing two maps along an
%          edge (« attachement de deux arêtes ») — the bi/bi, mon/bi and
%          mon/mon cases, with the effect on f (faces) and g (genus).
%          Page 54: pencil sketches only.
%
% Notation fixed for the batch. \mathfrak{P}_2(A) for his script P_2 (the
% set of two-element subsets); \smallsetminus for his set difference; A^* for
% his starred cores; \widehat{\Sigma} for his Σ with a hat. His « ens. »,
% « segm. », « comp. », « card » are kept as written. Brown-ink pages 48-50
% are a fast hand whose formulas carry and whose prose often does not; the
% rotated margins on pages 41-43 and 51 are largely \ill{}.
%
% The dating is the inventory's own for the folder, copied verbatim. Its
% group, [69-89] « Géométrie et topologie combinatoire ».
\documentclass[11pt,a4paper]{article}
% Le préambule commun à toutes les transcriptions du fonds Grothendieck.
%
% Il tient en une idée : **l'incertitude se compose, elle ne se résout pas.**
% Une transcription qui lisse un mot douteux a détruit la seule chose pour
% laquelle on la faisait — un lecteur doit pouvoir savoir, à chaque endroit,
% ce qui était sur le feuillet et ce qui a été ajouté. Les six macros qui
% suivent sont donc l'appareil critique, et non de la décoration.
%
% Les mêmes commandes sont reconnues par `scripts/render.mjs`, qui produit la
% vue de lecture du site. Une macro ajoutée ici doit l'être aussi là-bas,
% faute de quoi le rendu échouera bruyamment — ce qui est voulu : un rendu
% silencieusement faux est pire qu'un rendu absent.

\ProvidesPackage{grothendieck}[2026/08/08 Transcriptions du fonds Grothendieck]

\usepackage{amsmath,amssymb,amsthm}
\usepackage{mathrsfs}
\usepackage{stmaryrd}
% Les diagrammes commutatifs sont partout dans ce fonds : c'est la langue
% même de la géométrie algébrique de Grothendieck, et une transcription qui
% les rendrait en prose ne transcrirait rien.
\usepackage{tikz-cd}
% Français par défaut — c'est la langue des feuillets — mais l'anglais est
% chargé aussi : la traduction et le résumé sont des documents anglais, et la
% typographie française (espaces avant les deux-points) y serait une faute.
% Ces documents basculent par \selectlanguage{english} en tête de corps.
\usepackage[english,french]{babel}
\usepackage{fontspec}
\usepackage{geometry}
\geometry{a4paper,margin=25mm,marginparwidth=28mm}
\usepackage{marginnote}
\usepackage{xcolor}
% ulem, not soul. `soul`'s \st and \ul are built on a character-by-character
% scan that XeLaTeX's font machinery breaks: the document fails with an
% undefined control sequence rather than degrading. ulem does the same two jobs
% and is Unicode-engine safe. `normalem` stops it hijacking \emph, which the
% transcriptions use for Grothendieck's own emphasis.
\usepackage[normalem]{ulem}
\usepackage{hyperref}
\hypersetup{colorlinks=true,linkcolor=black,urlcolor=black,pdfborder={0 0 0}}

% --- Métadonnées du lot -----------------------------------------------------
% Reprises telles quelles par le script de rendu, qui les affiche en tête de
% la vue de lecture. Elles fixent ce que le fichier prétend couvrir : sans
% elles, une transcription flotte sans point d'attache dans un fonds de seize
% mille feuillets.

\newcommand{\folder}[1]{\gdef\grothFolder{#1}}
\newcommand{\batch}[1]{\gdef\grothBatch{#1}}
\newcommand{\leaves}[2]{\gdef\grothFirst{#1}\gdef\grothLast{#2}}
\newcommand{\foldertitle}[1]{\gdef\grothFoldertitle{#1}}
\gdef\grothFolder{?}\gdef\grothBatch{?}\gdef\grothFirst{?}\gdef\grothLast{?}\gdef\grothFoldertitle{}

% --- Le résumé --------------------------------------------------------------
%
% Il ouvre la lecture modernisée, sous le titre « Résumé », et il est composé
% en petit. Ce n'est pas un ornement : c'est le seul passage du document qui
% ne soit pas des mathématiques, et un lecteur doit pouvoir le reconnaître
% avant de l'avoir lu — puis le sauter s'il connaît déjà le sujet, ou s'y
% arrêter s'il ne le connaît pas. Le filet qui le suit marque la reprise du
% corps.

\newenvironment{resume}
  {\small\setlength{\parskip}{0.45em}}
  {\par\medskip\hrule\medskip}

% --- Mots-clés ---------------------------------------------------------------
%
% En anglais, en clôture du résumé de la lecture modernisée : le vocabulaire
% moderne sous lequel chercher ce que les feuillets construisent. Le manifeste
% du site (`npm run manifest`) les extrait pour étiqueter la cote entière —
% cette ligne est la source unique des tags, il n'y a pas de fichier à côté.
\newcommand{\keywords}[1]{\par\smallskip\noindent
  {\footnotesize\textsc{Keywords} --- \textit{#1}}\par}

% --- L'appareil critique ----------------------------------------------------

% Le numéro de feuillet, en marge. C'est l'ancre qui permet de régler un
% désaccord sur une formule en regardant *un* feuillet plutôt qu'une chemise
% de six cents. Le site s'en sert aussi pour tourner les pages du fac-similé
% quand on fait défiler la transcription.
\newcommand{\leaf}[1]{%
  \marginnote{\footnotesize\color{gray!70}\textsf{#1}}%
  \label{leaf:#1}\ignorespaces}

% Illisible : on ne devine pas. Un mot inventé qui se lit comme les autres est
% le pire résultat possible.
\newcommand{\ill}{\textcolor{red!60!black}{[\,\ldots\,]}}

% Lecture douteuse : le mot est donné, et signalé comme incertain.
\newcommand{\uncertain}[1]{\uline{#1}}

% Ajout de l'éditeur — une lettre manquante, un mot sous-entendu.
\newcommand{\add}[1]{\textcolor{blue!50!black}{[#1]}}

% Biffé par Grothendieck lui-même. Ce qu'il a rayé fait partie du document :
% une rature dit souvent où la pensée a changé de direction.
\newcommand{\struck}[1]{\sout{#1}}

% Note du transcripteur, sur l'état matériel du feuillet.
\newcommand{\note}[1]{\par\smallskip{\small\color{gray!60!black}\textit{[#1]}}\par\smallskip}

% Note marginale de Grothendieck, distincte du corps du texte.
\newcommand{\marginal}[1]{\par\smallskip\noindent\hspace{1em}%
  {\small\color{gray!60!black}#1}\par\smallskip}

% --- Titre ------------------------------------------------------------------

\AtBeginDocument{%
  \begin{center}
    {\large\bfseries Fonds Alexandre Grothendieck --- cote n\textsuperscript{o}~\grothFolder}\\[2pt]
    {\itshape\grothFoldertitle}\\[4pt]
    {\small feuillets \grothFirst--\grothLast\ (lot \grothBatch)}\\[2pt]
    {\footnotesize\color{gray!60!black}Transcription établie d'après le fac-similé numérisé
     par l'Université de Montpellier.\\
     Les lectures incertaines sont soulignées, les illisibles notées [\,\ldots\,],
     les ajouts entre crochets.}
  \end{center}
  \vspace{1em}}

\endinput


\folder{81}
\batch{3}
\pages{41}{54}
\dating{[à partir de 1981]}
\watermark{Édition de démonstration}
\foldertitle{Immersions du disque et de la circonférence : notes manuscrites (s.d.).}

\begin{document}

\section{Relations d'équi.\ non parallèles}
\note{ces mots sont de sa main, au crayon, sur le feuillet 40 (dernier du lot
précédent), par ailleurs vierge, qui ne reçoit pas de numéro de page}

\page{41}
Soit $S$ un ens.\ fini, \add{de} \uncertain{cardinal} $N$, avec une structure
polygonale ou bi-circulation $\{u,u^{-1}\}$ ($\operatorname{card} S\geqslant
3$).

\struck{\ill{}} $T\subset S$ \add{de cardinal $i$} est dit un
\struck{\ill{}} \emph{segment} de $S$, si
a) \struck{\ill{}} $T\neq S$, $T\neq\emptyset$ \struck{$i\neq N$} i.e.\
$0<i<N$
b) $\exists\, s_{0}\in S$, \struck{$i\in\mathbb{N}$}, tel que
$T=\{s_{0},us_{0},\dots,u^{i-1}s_{0}\}$ (alors, posant
$t_{0}=u^{i-1}s_{0}$, on a aussi $T=\{t_{0},vt_{0},\dots,v^{i-1}t_{0}\}$,
où $v=u^{-1}$).

\emph{Ex.} les parties réduites à un point. \struck{\ill{}} les
parties-arêtes.

$T$ est un segm.\ ssi $S\smallsetminus T$ en est un.

\emph{Intersection} de deux segments \add{contenus dans un segment} (en est un
ssi il est $\neq\emptyset$)

Réunion de deux segments qui se rencontrent et tels que $T\cup T'\neq S$
\add{est un segment}

Extrémités d'un segment : $s_{0}$, $t_{0}$ (distinctes ssi $i\neq1$ i.e.\ $T$
non réduit à un point).
\marginal{\uncertain{Propr.} : \ill{}, \uncertain{soit} $A$ \ill{}}
\marginal{i.e.\ \ill{} de $S$ \ill{}}

Origine, extrémité pour une orientation.

[\emph{NB} Si $u$ i.e.\ l'orientation choisi, le $s_{0}$ est unique i.e.\ il
y a corr.\ 1-1 entre l'ens.\ des segments de $S$ et l'ens.\
$(S,[1,N-1]_{\mathbb{N}})$, via
\[
  (s_{0},i)\longmapsto\{s_{0},us_{0},\dots,u^{i-1}s_{0}\}.
\]
\note{le crochet ouvert devant \emph{NB} n'est pas refermé sur la page ; le
couple $(S,[1,N-1]_{\mathbb{N}})$ est écrit tel quel, pour ce qui est
visiblement le produit $S\times[1,N-1]$}

Soit $S'\subsetneq S$ une partie de $S$. Alors \struck{les segments de $S$
contenus} deux segments de $S$ contenus dans $S'$, et qui se rencontrent, ont
comme réunion un segment \add{de $S$} contenu dans $S'$ — cela implique

\page{42}
que pour tout segment $I$ de $S$, \add{$I\subset S'$}, il existe un plus grand
segment de $S$ contenu dans $S'$ le contenant. Les éléments maximaux de
\add{$\operatorname{Segm}_{S'}(S)$} sont \uncertain{exactement} les
\emph{segments} \ill{} \ill{} partition de $S'$, \ill{} d'une relation
d'équivalence dans $S'$.
\marginal{i.e.\ \uncertain{introduisant}
$\operatorname{Segm}_{S'}(S)=\{I\subset S'\mid
I\in\operatorname{Segm}(S)\}$}
\marginal{les éléments maximaux de $\operatorname{Segm}(S,S')$ \ill{} les
segments \uncertain{composants} de $S'$ \ill{} en corr.\ \ill{}
$S\smallsetminus I$ \ill{}}

\ill{} d'une réalisation top.\ de $(S,\{u,u^{-1}\})$ par un polygone top.\
$(X,S)$, si $S''=S\smallsetminus S'$, les \struck{classes} segments de $S'$
sont en corr.\ biunivoque \add{avec les} comp.\ conn.\ de $X\smallsetminus
S''$ qui rencontrent $S'$ — ce sont les intersections de $S'$ avec ces comp.\
connexes.

\note{dans la marge gauche, un cercle portant une vingtaine de points, les
uns pleins, les autres évidés ; légende en dessous : « $S''$
\uncertain{marqué} \ill{} »}

\emph{Proposition} Soient $A$, $B$ deux parties \struck{disjointes}
\struck{non vides} de $S$, \ill{} \struck{$A$ et $B$} $A,B\neq\emptyset,S$,
$A\cap B=\emptyset$.

Conditions équivalentes :

a) Il existe des segments $\overline{A}$, $\overline{B}$ de $S$, tels que
$\overline{A}\supset A$, $\overline{B}\supset B$,
$\overline{A}\cap\overline{B}=\emptyset$,

\marginal{\emph{NB} On \ill{} \ill{} $\overline{A}$, $\overline{B}$
\add{plus petits} \struck{\ill{}} segments \ill{}}

\page{43}
b) Pour $a,a'\in A$, $b,b'\in B$, avec $a\neq a'$, $b\neq b'$,
\struck{les \ill{}} $a$, $a'$ sont \struck{contigus} \emph{adjacents}
\add{l'un à l'autre} pour la structure bi-circulaire sur $\{a,a',b,b'\}$
induite par celle de $S$, \struck{i.e.\ \ill{}} ($\Longleftrightarrow$
$\{b,b'\}$ sont adjacents)

$\Longleftrightarrow$ $\{a,a'\}$ ne \uncertain{sont} pas
\uncertain{opposés} $\Longleftrightarrow$ $b$, $b'$ ne \uncertain{sont}
\ill{}

c) $B$ est contenu dans un segment composant de $S\smallsetminus A$

c') $A$ est contenu dans un segment composant de $S\smallsetminus B$

\uncertain{Quand} il en est ainsi : Parmi les \add{couples de} segments
$(\overline{A},\overline{B})$ satisfaisant les conditions de a), il existe
un plus petit…

On dit que $A$, $B$ sont \emph{non parallèles} (plus généralement si
$A\cap B$ est un segment). \emph{Ex.} Si \struck{\ill{}} $A$ ou $B$ réduit
à un pt.\ (de cardinal 1), cela signifie \ill{} pour la bicirculation \ill{}
un côté du \uncertain{couvert} $A\cup B$ \ill{} $B$ est un côté). Donc, $A$
et $B$ sont non parallèles \uncertain{ssi} $\forall A'\in\mathfrak{P}_{2}(A)$,
$B'\in\mathfrak{P}_{2}(B)$, $A'$ et $B'$ sont non parallèles.

\marginal{d) Quand $S$ est réalisé comme partie \add{\uncertain{convexe}} du
\uncertain{bord} d'un disque plan (dans un plan réel affine) : Soit $P_{A}$
la réunion des segments affines joignant deux à deux les points de $A$ (i.e.\
l'enveloppe convexe de $A$) \ill{} $\Longleftrightarrow$ $P_{A}\cap
P_{B}=\emptyset$ \ill{} \struck{\ill{} $A$, $B$ \ill{}} \ill{} conditions
\ill{} $A\cup B$ \ill{} $C\supset A\cup B$, \ill{} les composantes \ill{}}
\note{longue note marginale en travers, écrite en rotation dans la marge
gauche et en partie surchargée ; seuls les termes mathématiques ressortent}

\page{44}
\emph{Structure à étudier} : Partitions de $S$ en des
$(S_{\alpha})_{\alpha\in J}$ mutuellement non parallèles. (On les appelle
des \emph{partitions non parallèles})
\marginal{\emph{Ex.} partitions \ill{} (au moins deux) \ill{} de $S$}

\emph{Intuition} À \struck{\ill{}} un ens.\ bicirculé \add{de cardinal $N$,}
muni d'une partition non parallèle, correspond (\uncertain{essentiellement}
pour les \ill{}) un ens.\ bicirculé (d'ordre \add{pair} cardinal $N'\leqslant
N$), avec une partition \add{non parallèle} de degré 2 (classes de cardinal
2)…

\note{sous ce texte, un grand cercle tracé à la main, sans aucune marque ; le
reste de la page est vide}

\page{45}
Soit $\Sigma$ un ens.\ de segments de $S$, satisfaisant \add{aux} les
conditions suivantes :

($\sigma_{1}$) \struck{\ill{}} Si $A,B\in\Sigma$, et si $A\cap
B\neq\emptyset$, alors $A\subset B$ ou $B\subset A$.

\add{($\sigma_{2}$)} Si de plus $A\neq B$ et $A\subset B$, alors $A$ est
contenu dans l'intérieur de $B$.

Posons alors, pour $A\in\Sigma$ :
\[
  A^{*}=A\smallsetminus\bigcup_{\substack{B\in\Sigma\\ B\subset A,\ B\neq A}}B .
\]
On a (par $\sigma_{2}$) $A^{*}\supset\partial A$ (ens.\ des extrémités de
$A$), donc $A^{*}\neq\emptyset$. Je dis que $A\mapsto A^{*}$,
$\Sigma\to\mathfrak{P}(S)$ est injectif, i.e.\ que si $A,B\in\Sigma$, alors
$A^{*}=B^{*}\Rightarrow A=B$ \struck{\ill{}} $A\neq B$ \struck{distincts}
$\Rightarrow A^{*}\neq B^{*}$. \add{Comme \ill{} $A^{*}$, $B^{*}$ sont
$\neq\emptyset$, il suffit de voir que $A^{*}\cap B^{*}=\emptyset$.}

\struck{En effet}, si $A\cap B=\emptyset$, alors $A^{*}\cap B^{*}=\emptyset$,
donc \ill{} (puisque $A^{*}$, $B^{*}$ non vides). Si $A\cap B\neq\emptyset$,
\ill{} (par $\sigma_{1}$) $A\subset B$ ou $B\subset A$, soit $A\subset B$.
\struck{Mais $A\cap\partial B=\emptyset$ (\struck{donc} par $\sigma_{2}$)
$A^{*}\cap\partial B=\emptyset$, donc $A^{*}\not\supset B^{*}$ (comme
\ill{} $\partial B\subset B^{*}$) \ill{}}
\note{ce passage est encadré puis biffé}

Comme $A\neq B$, $A\cap B^{*}=\emptyset$ par définition de $B^{*}$,
\struck{Je dis} Ainsi $A^{*}$, $B^{*}$ \ill{} $A^{*}\subset A$ segment
$\subset S\smallsetminus B^{*}$.

\marginal{\emph{Prop} Si $A,B\in\Sigma$, $A\neq B\Rightarrow A^{*}\cap
B^{*}=\emptyset$ (i.e.\ \ill{} $A^{*}\neq B^{*}$).}
\marginal{On a $\bigcup_{A\in\Sigma}A=\bigcup_{A\in\Sigma}A^{*}$. Si
$x\in\bigcup A$, alors $\Sigma(x)$ \ill{} non vide \ill{} $A\in\Sigma(x)$
minimal \ill{} $x\in A^{*}$ \ill{}}

\page{46}
\emph{Conséquences} : Si $\bigcup_{A\in\Sigma}A^{*}=S$, alors
\struck{les} $(A^{*})_{A\in\Sigma}$ est une partition \struck{dans}
\uncertain{par-dessus} de $S$. Si $\bigcup_{A\in\Sigma}A^{*}\neq S$, alors
posant \struck{\ill{}}
\[
  \begin{cases}
    R=S\smallsetminus\bigcup_{A\in\Sigma}A^{*}, &
    \widehat{\Sigma}=\Sigma\cup\{R\}\\[2pt]
    R^{*}=R\smallsetminus\bigcup_{A\in\widehat{\Sigma},\ A\subsetneq R}A=R &
  \end{cases}
\]
alors $(A^{*})_{A\in\widehat{\Sigma}}$ est une partition
\uncertain{par-dessus} de $S$.
\note{« par-dessus » est une lecture douteuse, ici comme deux lignes plus
haut ; la page 44 dit, pour la même chose, « partition non parallèle »}

\emph{NB} Si $R$ est lui-même un segment, alors \struck{\ill{}}
$\widehat{\Sigma}$ satisfait les conditions de $\Sigma$ — et $R$ est un
élément « isolé » de $\widehat{\Sigma}$, en ce sens que \add{pour} $A$
\ill{} élément de $\widehat{\Sigma}$, \struck{\ill{}} $R\cap A=\emptyset$ —
i.e.\ \ill{} minimal et \ill{}

Considérons

($\sigma_{3}$) $R=S\smallsetminus\bigcup_{A\in\Sigma}A$
($=S\smallsetminus\bigcup A^{*}$) n'est pas un segment

(\emph{NB} précisément ce cas si $R=\emptyset$ par ex.)

\marginal{\uncertain{Seul}}
\uncertain{Moyennant} donc ce cas \struck{\ill{}} $\Sigma$ est \add{\ill{}}
\struck{défini} \ill{} connaît la partition
$\widehat{\Sigma}^{*}=\{A^{*}\mid A\in\widehat{\Sigma}\}$
\[
  \widehat{\Sigma}=
  \begin{cases}
    \Sigma & \text{si } R=\emptyset\\
    \Sigma\cup\{R\} & \text{si } R\neq\emptyset
  \end{cases}
\]

\page{47}
Je veux construire $A$ en termes de $\widehat{\Sigma}^{*}$,
$A^{*}\in\widehat{\Sigma}^{*}$ (dans le cas $A\in\Sigma$). Bien sûr, $A$ est
un segment qui contient $A^{*}$, et comme $A^{*}\supset\partial A$, il est
minimal (\ill{}) parmi les segments qui contiennent $A^{*}$ — mais cela ne
suffit pas à le caractériser. Il a cependant la propriété supplémentaire,
$\forall B^{*}\in\widehat{\Sigma}^{*}$, ou $B^{*}\cap A=\emptyset$, ou
$B^{*}\subset A$.

\emph{NB} En fait, avec les seules conditions $\sigma_{1}$ à $\sigma_{3}$ sur
$\Sigma$, il n'est pas vrai que la connaissance de $\widehat{\Sigma}^{*}$
redonne $\Sigma$, exemple la relation d'équivalence $R$ \add{\ill{}} avec
trois classes $\{a_{1},a_{1}'\}$, $\{b,b'\}$, $\{a_{2},a_{2}'\}$ ($\Sigma_{1}$
formé des segments $\{a_{i},a_{i}'\}$ ($i\in\{1,2\}$), et
$\{a_{1},a_{1}',b,b'\}$, et $\Sigma_{2}$ défini de façon symétrique,
\ill{} la même $R$).

\note{à droite, un cercle traversé de trois cordes verticales parallèles,
d'extrémités notées $a_{1}$, $b$, $a_{2}$ en haut et $a_{1}'$, $b'$,
$a_{2}'$ en bas}

\struck{\emph{Exemple} plus général : on part \ill{} de façon unique \ill{}
donnons $\Sigma$ de façon \ill{} $\widehat{\Sigma}^{*}=R$, et
$R\overset{\mathrm{def}}{=}S\smallsetminus\bigsqcup_{A\in\Sigma}A=S_{i}$,
chaque donnée d'une classe de $R$ si non un segment \ill{}}
\note{ce dernier paragraphe est barré de deux traits en diagonale}

\page{48}
\emph{Relation d'équivalence non \uncertain{parallèle}} sur $S$.
\note{à partir d'ici, et jusqu'à la page 50, une encre brune et une main
rapide}

Segments \add{$R$-}\emph{saturés} : \uncertain{saturés} par $R$ \ill{} par
passage au complémentaire, et par réunion et intersection de segments
\add{p.ex.\ si} ($A\cup B\neq S$, $A\cap B\neq\emptyset$).

Si $S_{i}\in R$, alors les segments composants de $S\smallsetminus S_{i}$
sont saturés, donc \ill{} leurs complémentaires, i.e.\ les segments
\uncertain{minimaux} contenant $S_{i}$ \ill{} \add{segments} (\ill{}
\uncertain{composants} de $S_{i}$). \struck{Plus généralement}
\struck{Si $T\subset S$ est une partie de $S$ $R$-saturée, alors les
segments composants de}

\emph{Prop} Les segments saturés minimaux sont les $S_{i}\in R$ qui sont des
segments.
\marginal{\ill{} \uncertain{anglets} \ill{}}

Il est évident que si $S_{i}\in R$ est un segment, alors c'est un segment
saturé, et il est minimal (car \ill{} non vides). Inversement, soit $A$
segment saturé minimal, \struck{\ill{}} soit $S_{i}\in R$, $S_{i}\subset A$,
\struck{soit} Je dis que $S_{i}=A$, \struck{\ill{}} i.e.\ qu'il n'y a pas deux
$S_{i},S_{j}\in R$ distincts, $\subset A$. En effet, si $A_{i}$ ($A_{j}$) est
le segment \ill{} \uncertain{composant} de \ill{} qui contient $S_{i}$
($S_{j}$) \ill{} $\subset A$, alors la condition que $S_{i}$, $S_{j}$ sont
\ill{} se traduit par la condition \struck{\ill{}} que $A_{i}\cap
A_{j}=\emptyset$, donc $A_{i}$, $A_{j}$ sont des segments saturés $\subset A$
et $\neq A$, absurde.

\note{dans la marge gauche, un arc de cercle aux deux extrémités marquées
d'un point, la moitié supérieure repassée à l'encre}
\marginal{\uncertain{unicité} des \ill{} de la condition \ill{}}

\page{49}
\emph{Cor} Tout \struck{\ill{}} segment saturé contient un
« \uncertain{anglet} » — donc si $R$ n'est pas la relation grossière, il
existe au moins deux anglets (car il existe alors au moins deux segments
saturés, complémentaires l'un de l'autre).

\add{$\Sigma_{0}$ l'ens.\ des anglets de $R$}
\note{« anglet » est la lecture retenue pour un mot écrit ici six fois ; il
désigne une classe de $R$ qui est un segment}

Soit \ill{} la réunion des anglets, $S_{1}$ son complémentaire. Alors, si
$S_{1}\neq\emptyset$, $S_{1}$ a au moins deux composantes
($\operatorname{card}\pi_{0}(S_{1})\geqslant2$), sinon, $R$, \uncertain{serait}
un segment contenant deux anglets, absurde. Si $\operatorname{card}R_{1}=2$,
alors $R_{1}\in R$ \ill{}

\struck{\ill{}} de $R$ consiste en les anglets et en $R_{1}$ (car les
$S_{i}\subset R_{1}$ sont de card.\ $\geqslant2$, \struck{\ill{}} \ill{}
1 seulement aux anglets…)

Si $\operatorname{card}R_{1}\geqslant3$, $R_{1}$ est lui-même un
\uncertain{polygone}, et $R_{1}=R\mid R_{1}$ est une relation d'équivalence
non \uncertain{parallèle}. De plus, \struck{\ill{}} $S_{i}\subset R_{1}$ n'est
contenu dans un segment composant de $R_{1}$ — sinon il serait \ill{}
segment saturé contenu dans \ill{} \ill{} \ill{}), lequel contiendrait un
anglet $\subset R$. — absurde. Si $R_{1}$ est la relation grossière, on a
fini — $R$ est formé des anglets, et de $R_{1}$.
\note{sur cette page, $R_{1}$ désigne tantôt la relation induite sur
$S_{1}$, tantôt $S_{1}$ lui-même ; transcrit tel qu'écrit}

\page{50}
Sinon, considérons l'ens.\ des \add{$R_{1}$-}anglets \struck{\ill{}}
\add{\ill{}} — il y en a au moins deux. Ces anglets, \ill{} c'est un
\uncertain{cour.} 1-1 \ill{} un ens.\ de segments $\Sigma_{2}$ de $S$, qui
sont \ill{} contenant au moins un \ill{} de $S_{1}$, et dont les extrémités
sont \uncertain{distinctes} de $S_{1}$, et qui sont mutuellement disjoints.
Soit $T_{2}$ \ill{}

\[
  \begin{array}{ccccc}
    \Sigma_{1} & \Sigma_{2} & \Sigma_{3} & \cdots & \Sigma_{i}\\
    T_{1} & T_{2} & T_{3} & & T_{i}
  \end{array}
\]

$\Sigma_{i}$ est un ens.\ de segments ($R$-saturés), mutuellement disjoints
\uncertain{(\ill{})}, $\operatorname{card}\Sigma_{i}\geqslant2$,
\add{de réunion $T_{i}$} (\emph{NB} les segments composants de $T_{i}$
\ill{} réunions partielles de $\Sigma_{i}$, \ill{})
\marginal{$i\geqslant1$}
\marginal{$i\geqslant2$ \ill{}}
\struck{tout élément $A\in\Sigma_{i-1}$ est contenu dans un
$B\in\Sigma_{i}$, tel \ill{}} Tout $B\in\Sigma_{i}$ contient au moins un
\add{segment composant} \struck{$B\in\Sigma_{i-1}$} de $T_{i-1}$, et
$\partial B\subset S\smallsetminus T_{i-1}$.

$\Sigma_{1}$ ens.\ de segments mutuellement disj.\ de $S=S_{0}$,
\struck{$\operatorname{card}(\Sigma_{1})\geqslant2$}
\[
  T_{1}=\bigcup_{A\in\Sigma_{1}}A,\qquad S_{1}=S\smallsetminus T_{1}
\]
\struck{$\operatorname{card}S_{1}\geqslant2$}
les segm.\ comp.\ de $S_{1}$ dans $S$ \uncertain{ont} au moins 2 \ill{}

$\Sigma_{2}'$ ens.\ de segments disjoints de $S_{1}$ \emph{non contenus}
dans un des segments composants \ill{},
$\operatorname{card}\Sigma_{1}\geqslant2$,
\[
  T_{2}=T_{1}\cup\bigcup_{A\in\Sigma_{2}'}A,\qquad
  S_{2}=S\smallsetminus T_{2},\qquad
  \operatorname{card}\pi_{0}(S_{2})\neq1
\]

\page{51}
$I$ \uncertain{quasi-segment} (\uncertain{fini}) ayant au moins deux
\uncertain{sommets}
\note{à partir d'ici une encre noire ; reprise, pour un « quasi-segment »
$I$, des notions des pages 41 à 47}

Deux \add{parties} d'éléments $\varepsilon,\varepsilon'\in\mathfrak{P}_{2}(I)$
sont dites \emph{non parallèles} si
\add{$\varepsilon\cap\varepsilon'=\emptyset$ et si}
$I_{\varepsilon}\cap I_{\varepsilon'}=\emptyset$ \add{ou
$I_{\varepsilon}\subset I_{\varepsilon'}$ ou $I_{\varepsilon'}\subset
I_{\varepsilon}$} (où $I_{\varepsilon}$ est le segment déterminé par un ens.\
à deux éléments $\varepsilon$).

Deux parties $A$, $A'$ de $I$, sont dites non parallèles si

a) $A\cap A'=\emptyset$

b) $\forall\varepsilon\in\mathfrak{P}_{2}(A)$,
$\varepsilon'\in\mathfrak{P}_{2}(A')$, $\varepsilon$, $\varepsilon'$ non
parallèles.

\struck{\ill{}} Cela revient à dire que $A'$ est contenu dans l'un des
\struck{\ill{}} composants de $I\smallsetminus A$ ou (si $A$ ne contient pas
d'extrémité de $I$) dans la réunion des deux composants extrêmes (i.e.\ ceux
qui contiennent les deux extrémités).
\marginal{\ill{} ou encore, que l'un des deux $A$, $A'$ est contenu dans un
segment \ill{} de $I$ qui ne \uncertain{rencontre} pas l'autre \ill{}}
\marginal{\emph{NB} un segment \ill{} une seule extrémité \ill{}}

Partitions non parallèles de $I$

$\Sigma$ ens.\ de segments de $I$ tel que :

a) Si $A,B\in\Sigma$, ou $A\cap B=\emptyset$, ou $A\subset B$, ou $B\subset A$

b) Si $A,B\in\Sigma$, $B\subsetneq A$, alors $B\cap\partial A=\emptyset$
\note{le signe entre $B$ et $\partial A$ est surchargé ; la lecture
$B\cap\partial A=\emptyset$ est celle que demande la condition ($\sigma_{2}$)
de la page 45}

c) \struck{\ill{}} $\partial I\subset\bigcup_{A\in\Sigma}A$ donc $\partial
I\subset\bigcup_{A\in\Sigma}\partial A$

Pour un tel $\Sigma$, posons pour $A\in\Sigma$
\[
  A^{*}=A\smallsetminus\bigcup_{\substack{B\in\Sigma\\ B\subsetneq A}}B
  \qquad(\supset\partial A\text{, donc }A^{*}\neq\emptyset)
\]
alors \struck{les $A^{*}$ sont deux à deux} pour $A,B\in\Sigma$, $A\neq B$,
on a $A^{*}$ et $B^{*}$ non parallèles — posons \struck{\ill{}} alors
\[
  I'=\bigcup_{A\in\Sigma}A=\bigcup_{A\in\Sigma}A^{*},\qquad
  R=I\smallsetminus I',\qquad R^{*}=R
\]
\[
  \widehat{\Sigma}=
  \begin{cases}
    \Sigma & \text{si } R=\emptyset\\
    \Sigma\cup\{R\} & \text{si } R\neq\emptyset
  \end{cases}
\]

\page{52}
alors
\[
  \widehat{\Sigma}^{*}=\{A^{*}\mid A\in\widehat{\Sigma}\}
\]
est une partition non parallèle de $I$.

\emph{NB} Toute partition non parallèle peut s'obtenir ainsi, mais de
plusieurs façons…

Soit $\Sigma_{\max}$ l'ens.\ des éléments maximaux de $\Sigma$. On voit de
suite que deux élts distincts de $\Sigma_{\max}$ sont disjoints, et que
\[
  I'=\bigcup_{A\in\Sigma_{\max}}A
\]
\note{la page s'arrête ici ; le reste en est vide}

\page{53}
\note{feuille carrée, d'un autre sujet : le recollement de deux cartes le long
d'une arête. En tête, trois croquis au crayon, repassés à l'encre aux
points de recollement : deux lobes accolés marqués « bi » et « bi », joints
par un trait vertical dont les extrémités portent des flèches courbes ; deux
lobes « mon » et « bi » séparés par une bande à deux traits parallèles, avec
une double flèche en haut et en bas ; deux lobes « mon » et « mon », une
double flèche au-dessus. En bas à droite, un cercle au crayon avec des
cordes inscrites}

Attachement de deux arêtes \struck{\uncertain{bipartites}} : on trouve
\add{paire d'}\struck{deux} arêtes \add{$a$, $a'$} \emph{binaires}
\struck{opposées} (qui sont telles que $(a,\omega)$, $(a',\omega)$ et
$(a,\overline{\omega})$, $(a',\overline{\omega})$ sont des \add{paires d'}
\struck{paires} associés, \struck{mais}
\[
  f=f'+f''-2,\qquad g=g'+g''
\]
\marginal{4 faces deviennent 2}

Attachement de deux arêtes, l'une monaire, l'autre binaire. On trouve une
paire d'arêtes $\{a,a'\}$ (\add{monaires} \struck{binaires}) incidentes :
\ill{} \ill{} faces
\[
  f^{*}=f'+f''-2,\qquad g=g'+g''
\]
\marginal{trois faces deviennent 1}

Attachement des deux arêtes \struck{\uncertain{monaires}}, on trouve
\struck{deux} paire d'arêtes monaires, incidentes à des faces distinctes
\[
  f=f'+f'',\qquad g=g'+g''-1
\]
\marginal{deux faces deviennent 2}
\note{chaque bilan est écrit à droite des formules, séparé d'elles par un
trait vertical ; le $g$ de la dernière ligne est souligné. Dans le deuxième
cas il écrit $f^{*}$ là où les deux autres ont $f$}

\page{54}
\note{page sans texte : croquis au crayon d'une même opération de collage
en volume — un disque vu en perspective relié par un tube court à une
feuille rectangulaire ; un petit cercle traversé d'un trait entre deux
arcs ; un grand rectangle à qui une bande est attachée par le haut ; une
feuille en perspective portant deux petits cylindres dressés. En haut à
droite, à l'encre, le trait vertical à extrémités recourbées du premier
croquis de la page 53}

\end{document}
